\documentclass[11pt,a4paper]{article}
\usepackage[utf8]{inputenc}
\usepackage{amsmath,amssymb,amsthm}
\usepackage{listings}
\usepackage{geometry}
\usepackage{hyperref}
\geometry{margin=2.5cm}

\newtheorem{theorem}{Theorem}
\newtheorem{lemma}[theorem]{Lemma}

\lstset{
  language=C++,
  basicstyle=\ttfamily\small,
  keywordstyle=\bfseries,
  breaklines=true,
  frame=single,
  numbers=left,
  numberstyle=\tiny,
  tabsize=4,
  showstringspaces=false
}

\title{IOI 2008 -- Fish}
\author{Solution}
\date{}

\begin{document}
\maketitle

\section{Problem Statement}
In a lake there are $N$ fish of distinct sizes, each belonging to one of $K$ species. A subset $S$ is \textbf{valid} if, for every species~$t$, the largest and smallest fish of species~$t$ in $S$ satisfy $\max \le 2 \cdot \min$. Count the number of non-empty valid subsets modulo a prime~$p$.

\section{Solution}

\subsection{Independence Across Species}
\begin{lemma}
The validity constraint is independent across species: for each species, we require $\max / \min \le 2$ within that species, regardless of other species. Therefore:
\[
\text{(valid subsets including empty)} = \prod_{t=1}^{K} (\text{valid subsets of species } t \text{, including empty}).
\]
\end{lemma}

\subsection{Counting Per-Species Valid Subsets}
For species~$t$ with fish sizes $s_1 < s_2 < \cdots < s_m$ (sorted), a valid subset is any subset $\{s_i, \ldots\}$ with $\max \le 2 \cdot \min$. We partition by the choice of minimum element.

For minimum element $s_i$, the eligible elements are those with size in $[s_i, 2s_i]$. Let $f(i) = |\{j > i : s_j \le 2s_i\}|$ (elements strictly larger than $s_i$ but within the $2\times$ bound). The number of subsets with minimum exactly $s_i$ is $2^{f(i)}$ (include $s_i$ and choose any subset of the $f(i)$ eligible larger elements).

Using two pointers (since $f(i)$ is non-decreasing in $i$), we compute $\sum_i 2^{f(i)}$ for each species in total $O(m)$ time.

\subsection{Final Answer}
\[
\text{Answer} = \left(\prod_{t=1}^{K} \Bigl(1 + \sum_{i} 2^{f_t(i)}\Bigr)\right) - 1,
\]
where the $-1$ removes the global empty subset.

\subsection{Complexity}
\begin{itemize}
    \item \textbf{Time:} $O(N \log N)$ for sorting, $O(N)$ for the two-pointer computation.
    \item \textbf{Space:} $O(N)$.
\end{itemize}

\section{C++ Solution}

\begin{lstlisting}
#include <bits/stdc++.h>
using namespace std;

long long power(long long base, long long exp, long long mod) {
    long long result = 1;
    base %= mod;
    while (exp > 0) {
        if (exp & 1) result = result * base % mod;
        base = base * base % mod;
        exp >>= 1;
    }
    return result;
}

int main(){
    ios::sync_with_stdio(false);
    cin.tie(nullptr);

    int N, K;
    long long MOD;
    cin >> N >> K >> MOD;

    vector<long long> sz(N);
    vector<int> sp(N);
    for (int i = 0; i < N; i++) {
        cin >> sz[i] >> sp[i];
        sp[i]--;
    }

    // Group fish by species, sorted by size
    vector<vector<long long>> bySpecies(K);
    for (int i = 0; i < N; i++)
        bySpecies[sp[i]].push_back(sz[i]);
    for (int t = 0; t < K; t++)
        sort(bySpecies[t].begin(), bySpecies[t].end());

    long long ans = 1;
    for (int t = 0; t < K; t++) {
        auto& sizes = bySpecies[t];
        int m = sizes.size();
        if (m == 0) continue;

        long long speciesCount = 1; // empty subset of this species
        int j = 0;
        for (int i = 0; i < m; i++) {
            while (j < m && sizes[j] <= 2 * sizes[i]) j++;
            int f = j - i - 1; // elements in (sizes[i], 2*sizes[i]]
            speciesCount = (speciesCount + power(2, f, MOD)) % MOD;
        }
        ans = ans * speciesCount % MOD;
    }

    ans = (ans - 1 + MOD) % MOD; // subtract the global empty subset
    cout << ans << "\n";
    return 0;
}
\end{lstlisting}

\section{Notes}
The per-species independence is the key structural insight. Within each species, the two-pointer technique efficiently counts valid subsets by tracking the rightmost element within the $2\times$ ratio bound. The pointer $j$ is monotonically non-decreasing, giving amortized $O(1)$ per element.

\end{document}
